% ID: SEL-D0068-OB-E03
% ESTADO: cotejada. Segunda pasada visual y borrador compilado revisados.
% Fuente: Selectividad/sel_2009_matematicas_aplicadas/mat apli CC SS02.pdf
% SHA-256: 24ee6d89c206838cb0cf4f32c4e54391e7ca81f77c4dd0be651ad7218f37beb6
% Página: 2; opción B; ejercicio original 3.

\noindent\underline{Parte I}

En una editorial hay dos máquinas A y B que encuadernan 100 y 900 libros al día, respectivamente. Además, se sabe que la probabilidad de que un libro encuadernado por A tenga algún fallo de encuadernación es del $2\%$, y del $10\%$ si ha sido encuadernado por la máquina B. Se elige, al azar, un libro encuadernado por esa editorial.
\begin{enumerate}[label=\alph*)]
\item \textbf{(1 punto)} Calcule la probabilidad de que no sea defectuoso.
\item \textbf{(1 punto)} Si es defectuoso, halle la probabilidad de haber sido encuadernado por la máquina A.
\end{enumerate}
\begin{samepage}
\noindent\underline{Parte II}

El tiempo que se tarda en la caja de un supermercado en cobrar a los clientes sigue una ley Normal con media desconocida y desviación típica $0.5$ minutos. Para una muestra aleatoria de 25 clientes se obtuvo un tiempo medio de $5.2$ minutos.
\begin{enumerate}[label=\alph*)]
\item \textbf{(1 punto)} Calcule un intervalo de confianza, al nivel del $97\%$, para el tiempo medio que se tarda en cobrar a los clientes.
\item \textbf{(1 punto)} Indique el tamaño muestral mínimo necesario para estimar dicho tiempo medio con un error máximo de $0.5$ y un nivel de confianza del $96\%$.
\end{enumerate}
\end{samepage}
